% Problem-section file following ../_template.tex.
\section{Convergence of the JRF iteration for mixed-state discrimination}
\label{sec:problem-63}

\paragraph{Problem.}
Does the Je\v{z}ek--\v{R}eh\'{a}\v{c}ek--Fiur\'{a}\v{s}ek (JRF) iteration,
initialized by the uniform POVM, converge to a globally optimal
minimum-error measurement for every finite ensemble containing mixed quantum
states?  Let $m\geq2$, let $p_i>0$ be prior probabilities, and let $\rho_i$ be
density operators on a finite-dimensional Hilbert space.  Define
\begin{equation}
  \xi_i:=p_i\rho_i,
  \qquad
  \operatorname{Tr}\rho_i=1,
  \qquad
  \sum_{i=1}^{m}p_i=1,
  \qquad
  \mathcal H_0:=\operatorname{supp}\!\left(\sum_{i=1}^{m}\xi_i\right),
  \label{eq:p63-weighted-ensemble}
\end{equation}
where at least one $\rho_i$ has rank greater than one.  On the signal space
$\mathcal H_0$ in Eq.~\eqref{eq:p63-weighted-ensemble}, the optimal guessing
probability is
\begin{equation}
  P_{\mathrm{opt}}
  :=\max_{\substack{\Pi_i\succeq0\\
                    \sum_{i=1}^{m}\Pi_i=I_{\mathcal H_0}}}
       \sum_{i=1}^{m}\operatorname{Tr}(\xi_i\Pi_i).
  \label{eq:p63-optimal-guessing-probability}
\end{equation}
Starting from $\Pi_i^{(0)}:=I_{\mathcal H_0}/m$, define the JRF iterates by
\begin{equation}
  \Lambda_k
  :=\left(\sum_{j=1}^{m}
       \xi_j\Pi_j^{(k)}\xi_j\right)^{1/2},
  \qquad
  \Pi_i^{(k+1)}
  :=\Lambda_k^{-1}\xi_i\Pi_i^{(k)}\xi_i\Lambda_k^{-1}
  \quad (i=1,\ldots,m).
  \label{eq:p63-jrf-iteration}
\end{equation}
For this initialization, $\Lambda_k$ is inverted on $\mathcal H_0$ and the
operators in Eq.~\eqref{eq:p63-jrf-iteration} form a POVM at every step.
Determine whether there always exists an optimal POVM
$\{\Pi_i^\star\}_{i=1}^{m}$ attaining
Eq.~\eqref{eq:p63-optimal-guessing-probability} such that
\begin{equation}
  \lim_{k\to\infty}
       \sum_{i=1}^{m}\lVert\Pi_i^{(k)}-\Pi_i^\star\rVert_2=0,
  \qquad
  \lim_{k\to\infty}
       \sum_{i=1}^{m}\operatorname{Tr}(\xi_i\Pi_i^{(k)})
       =P_{\mathrm{opt}},
  \label{eq:p63-global-convergence}
\end{equation}
where $\lVert\cdot\rVert_2$ is the Hilbert--Schmidt norm.  If
Eq.~\eqref{eq:p63-global-convergence} fails, construct an explicit
mixed-state counterexample and characterize its limiting behavior.

\subsection*{Status}

Unsolved

\subsection*{Source}

Je\v{z}ek, \v{R}eh\'{a}\v{c}ek, and Fiur\'{a}\v{s}ek introduced the
iteration and explicitly reported that its numerically observed global
convergence lacked a general proof
\sourcecite{ref:p63-jezek-rehacek-fiurasek}{JRF02}.  L\"u and Dong proved
the pure-state case and explicitly left convergence for general mixed-state
ensembles open \sourcecite{ref:p63-lu-dong}{LD26}.

\subsection*{Progress}

\begin{itemize}
  \item The original work observed monotonic convergence to the global optimum
  in extensive tests on ensembles of up to four pure or mixed states in
  dimensions two through four, but supplied neither a proof nor a
  counterexample for arbitrary ensembles
  \sourcecite{ref:p63-jezek-rehacek-fiurasek}{JRF02}.

  \item Tyson's directional-iterate framework proves the general monotonicity
  bound
  \begin{equation}
    P_k\leq\operatorname{Tr}\Lambda_k
        \leq P_{k+1}\leq P_{\mathrm{opt}},
    \qquad
    P_k:=\sum_{i=1}^{m}\operatorname{Tr}(\xi_i\Pi_i^{(k)}),
    \label{eq:p63-monotonicity}
  \end{equation}
  which L\"u and Dong rederive from a polar decomposition
  \sourcecite{ref:p63-tyson}{Tys10},
  \sourcecite{ref:p63-lu-dong}{LD26}.  Equation~\eqref{eq:p63-monotonicity}
  implies that $P_k$ has a limit, but does not show that the limit equals
  $P_{\mathrm{opt}}$ or that the POVM sequence converges.

  \item L\"u and Dong prove convergence to an optimal measurement for every
  pure-state ensemble under a mild support condition satisfied by the uniform
  initialization in Eq.~\eqref{eq:p63-jrf-iteration}.  They also obtain a
  conditional result for a special embedded class of mixed-state ensembles,
  but explain that their accumulation-point argument does not extend to
  arbitrary mixed states \sourcecite{ref:p63-lu-dong}{LD26}.
\end{itemize}

\subsection*{References}

\begin{enumerate}[label={},leftmargin=4.8em,labelsep=0.5em]
  \footnotesize
  \item[\textup{[JRF02]}]\label{ref:p63-jezek-rehacek-fiurasek}
  M. Je\v{z}ek, J. \v{R}eh\'{a}\v{c}ek, and J. Fiur\'{a}\v{s}ek,
  ``Finding Optimal Strategies for Minimum-Error Quantum-State
  Discrimination,'' \emph{Physical Review A} \textbf{65}, 060301(R) (2002).
  \href{https://doi.org/10.1103/PhysRevA.65.060301}%
       {doi:10.1103/PhysRevA.65.060301};
  \href{https://arxiv.org/abs/quant-ph/0201109}%
       {arXiv:quant-ph/0201109}.

  \item[\textup{[Tys10]}]\label{ref:p63-tyson}
  J. Tyson,
  ``Two-Sided Bounds on Minimum-Error Quantum Measurement, on the
  Reversibility of Quantum Dynamics, and on the Maximum Overlap Problem Using
  Directional Iterates,'' \emph{Journal of Mathematical Physics} \textbf{51},
  092204 (2010).
  \href{https://doi.org/10.1063/1.3463451}{doi:10.1063/1.3463451};
  \href{https://arxiv.org/abs/0907.3386}{arXiv:0907.3386}.

  \item[\textup{[LD26]}]\label{ref:p63-lu-dong}
  X. L\"u and S.-H. Dong,
  ``Iterative Algorithm for Minimum-Error Quantum State Discrimination:
  Convergence for Pure-State Ensembles,''
  \emph{Physical Review A} \textbf{113}, 022451 (2026).
  \href{https://doi.org/10.1103/q7wq-ygm9}{doi:10.1103/q7wq-ygm9}.
\end{enumerate}

\subsection*{Comment}

The pure-state case is solved.  The remaining issue is whether the uniform
initialization prevents nonoptimal accumulation points or nonconvergent
last-iterate behavior for arbitrary mixed-state ensembles.  The existence of
convergent semidefinite-programming methods for
Eq.~\eqref{eq:p63-optimal-guessing-probability} does not establish convergence
of the specific nonlinear map in Eq.~\eqref{eq:p63-jrf-iteration}.

\subsection*{Tag}

Quantum state discrimination; Quantum measurement theory; Convex optimization;
Semidefinite programming

\subsection*{ID}

\texttt{op\_76e284219621a785}
